% Fichier complété par tools/complete_missing_corrections.py.
% Source : CIN/CIN-02-VitesseAcceleration/1025_RTR/1025_RTR.tex
% Statut : rédigé (2 question(s)).
\begin{corrigebox}[Corrigé rédigé]
\CorrigeQuestion{1}
On applique la composition des vitesses puis la formule de changement de point :
                $\vec V(P,2/0)=\vec V(A,2/0)+\vec\Omega_{2/0}\wedge\overrightarrow{AP}$.
                Avec $\vec\Omega_{2/0}=\vec\Omega_{2/1}+\vec\Omega_{1/0}$, on obtient
                successivement les vitesses de $B$, $D$ et $C$ en remplaçant
                $\overrightarrow{AB}$, $\overrightarrow{AD}$ et $\overrightarrow{AC}$ par
                les vecteurs donnés sur la figure. Les composantes communes au solide 2
                diffèrent uniquement du terme $\vec\Omega_{2/0}\wedge\overrightarrow{BP}$.
\CorrigeQuestion{2}
Les accélérations sont obtenues par
                $\vec\Gamma(P,2/0)=\vec\Gamma(B,2/0)+\dot{\vec\Omega}_{2/0}\wedge
                \overrightarrow{BP}+\vec\Omega_{2/0}\wedge(\vec\Omega_{2/0}\wedge
                \overrightarrow{BP})$. Cette expression fournit directement les termes
                tangentiels (proportionnels aux accélérations angulaires) et normaux
                (proportionnels aux carrés des vitesses angulaires) pour $D$ et $C$.
\end{corrigebox}
